% ===== Appendix: corrections to the book, what we added, sources (after Questions; not presented) =====
\appendix
\partaccent{lavender}\gdef\idiottopicemoji{books}\gdef\idiottopicname{Հավելված}

\begin{frame}{Հավելված Ա. Ուղղումներ գրքի նկատմամբ}
  \scriptsize
  \setlength{\tabcolsep}{4pt}
  \renewcommand{\arraystretch}{1.25}
  \begin{tabular}{@{}L{2.1cm}L{3.0cm}L{4.8cm}L{3.6cm}@{}}
    \toprule
    \textbf{Թեմա} & \textbf{Գրքում (2016)} & \textbf{Այսօր} & \textbf{Աղբյուր}\\
    \midrule
    Սերոտոնին ու դեպրեսիա & Մոնոամինային վարկածը քննադատվում է, բայց կարևոր է & «Սերոտոնինի պակասի» համոզիչ ապացույց չկա (վիճելի ակնարկ)։ Հակադեպրեսանտները պլացեբոյից արդյունավետ են & Moncrieff et al., 2022; Jauhar et al., 2023; Cipriani et al., Lancet, 2018\\
    Կոկաին & «Նմանակում է դոֆամինը» & Փակում է դոֆամինի տրանսպորտերը (հետադարձ կլանումը) & NIDA, Cocaine Research Report\\
    Դոֆամին & Պարգևի ուղին հաճույք է տալիս & Դոֆամինը ավելի շատ «ուզել» է, քան «հավանել» & Berridge \& Robinson, Brain Res. Rev., 1998\\
    Ռոբին Ուիլյամս & Դեպրեսիայի ու «եսասեր» պիտակի օրինակ & Հետմահու՝ Լևիի մարմնիկներով դեմենցիա & Schneider Williams, Neurology, 2016\\
    «3-րդ դեմքի» ձայներ & Շիզոֆրենիայի հուսալի նշան & DSM-5-ը (2013) առանձին կշիռ չի տալիս. լինում են նաև այլ վիճակներում & Tandon et al., Schizophr. Res., 2013\\
    \bottomrule
  \end{tabular}
\end{frame}

\begin{frame}{Հավելված Բ. Ինչ ենք ավելացրել գրքից դուրս}
  \begin{columns}[T]
    \begin{column}{0.48\textwidth}
      \begin{card}[colback=mint!30!cardwhite, fontupper=\scriptsize]
        {\blackfont Թվեր ու հետազոտություններ}\par\smallskip
        \begin{itemize}\setlength{\itemsep}{0.5mm}
          \item Հոգեկան խանգարումներ՝ 29\% կյանքում, 18\% վերջին տարում՝ Steel et al., 2014
          \item Դեպրեսիայի ժառանգականություն՝ 37\%՝ Sullivan et al., 2000 (գրքի աղբյուրը)
          \item Հոլմս-Ռահե՝ ամուսնություն 50, Սուրբ Ծնունդ 12 (նույն սանդղակից)
          \item Ուղղումների աղբյուրները (Ա աղյուսակ)
        \end{itemize}
      \end{card}
    \end{column}
    \begin{column}{0.48\textwidth}
      \begin{card}[colback=peach!35!cardwhite, fontupper=\scriptsize]
        {\blackfont Մերն է, ոչ թե գրքից}\par\smallskip
        \begin{itemize}\setlength{\itemsep}{0.5mm}
          \item Սխեմատիկ գրաֆիկները (SSRI-ի ազդեցության ժամանակը, սթրեսի կորերը, Յերկս-Դոդսոն, տոլերանտության «սանդուղքը», ցավի ճոճանակը)՝ գրքի նկարագրությամբ, առանց իրական տվյալների
          \item «Փորձիր տանը» առաջարկները, շարքի ամփոփումը, վերջին խոսքը
          \item Մեմերը. բովանդակությունը գրքից է
        \end{itemize}
      \end{card}
    \end{column}
  \end{columns}
\end{frame}

\begin{frame}{Հավելված Գ. Աղբյուրներ. գլուխ 8 (1/2)}
  \scriptsize
  Գրքի ծանոթագրություններից, որոնց վրա հիմնված են սլայդները\par\smallskip
  \begin{columns}[T]
    \begin{column}{0.48\textwidth}
      \begin{itemize}\setlength{\itemsep}{0.3mm}
        \item Ramachandran V. S., Hubbard E. M. Synaesthesia. J. Conscious. Stud., 2001
        \item Hirschfeld R. History of the monoamine hypothesis of depression. J. Clin. Psychiatry, 2000
        \item Adrien J. Sleep and depression. Sleep Med. Rev., 2002
        \item Auer D. P. et al. Glutamate in the anterior cingulate in depression. Biol. Psychiatry, 2000
        \item Lok A. et al. HPA axis in recurrent depression. Psychoneuroendocrinology, 2012
        \item Eyre H., Baune B. T. Neuroplastic changes in depression. Psychoneuroendocrinology, 2012
        \item Katon W. et al. Depression and dementia risk. Arch. Gen. Psychiatry, 2012
        \item Epp A. M. et al. Stroop task in depression. Clin. Psychol. Rev., 2012
      \end{itemize}
    \end{column}
    \begin{column}{0.48\textwidth}
      \begin{itemize}\setlength{\itemsep}{0.3mm}
        \item Sullivan P. F. et al. Genetic epidemiology of major depression. Am. J. Psychiatry, 2000 (գրքում՝ 2007)
        \item Holmes T. H., Rahe R. H. The social readjustment rating scale. J. Psychosom. Res., 1967
        \item Barrett D. H. et al. Cognitive functioning and PTSD. Am. J. Psychiatry, 1996
        \item Broadhurst P. L. Emotionality and the Yerkes-Dodson law. J. Exp. Psychol., 1957
        \item Ulrich R. S. et al. Stress recovery in natural and urban environments. J. Environ. Psychol., 1991
        \item Dedovic K. et al. The brain and the stress axis. NeuroImage, 2009
        \item Monroe S. M., Harkness K. L. Life stress and the kindling hypothesis. Psychol. Rev., 2005
        \item Thoumi F. E. The numbers game. J. Drug Issues, 2005
      \end{itemize}
    \end{column}
  \end{columns}
\end{frame}

\begin{frame}{Հավելված Դ. Աղբյուրներ. գլուխ 8 (2/2)}
  \scriptsize
  Գրքի ծանոթագրություններից, որոնց վրա հիմնված են սլայդները\par\smallskip
  \begin{columns}[T]
    \begin{column}{0.48\textwidth}
      \begin{itemize}\setlength{\itemsep}{0.3mm}
        \item Caine S. B. et al. Cocaine self-administration in D3 knockout mice. Exp. Clin. Psychopharmacol., 2012
        \item Dalley J. W. et al. Impulse control and prefrontal serotonin. Neuropsychopharmacology, 2002
        \item Robinson T. E., Berridge K. C. Incentive-sensitization theory of addiction. Brain Res. Rev., 1993
        \item Brown R. Arousal and gambling addiction. Subst. Use Misuse, 1986
        \item Everitt B. J. et al. Amygdala and drug cues. Ann. N. Y. Acad. Sci., 1999
        \item Robinson G. M. et al. Patients who inject methadone syrup. Drug Alcohol Rev., 2000
        \item Clark L., Robbins T. W. Decision-making deficits in drug addiction. Trends Cogn. Sci., 2002
        \item Kreek M. J. et al. Genetics, impulsivity and addiction. Nat. Neurosci., 2005
      \end{itemize}
    \end{column}
    \begin{column}{0.48\textwidth}
      \begin{itemize}\setlength{\itemsep}{0.3mm}
        \item Shergill S. S. et al. Auditory hallucinations in schizophrenia. Am. J. Psychiatry, 2000
        \item Allen P. et al. The hallucinating brain. Neurosci. Biobehav. Rev., 2008
        \item Blakemore S.-J. et al. Self-produced stimuli and auditory hallucinations. Psychol. Med., 2000
        \item Buckner R. L., Carroll D. C. Self-projection and the brain. Trends Cogn. Sci., 2007
        \item Young A. W. et al. The Capgras and Cotard delusions. Psychopathology, 1994
        \item Coltheart M. et al. Delusional belief. Annu. Rev. Psychol., 2011
        \item Corlett P. et al. Toward a neurobiology of delusions. Prog. Neurobiol., 2010
        \item Coyle J. T. Glutamatergic dysfunction in schizophrenia. Harv. Rev. Psychiatry, 1996
      \end{itemize}
    \end{column}
  \end{columns}
\end{frame}
